\documentclass[
  lang=cn,
  color=blue,
  mode=fancy,
  scheme=chinese,
  chinesefont=nofont,
  math=newtx,
  12pt
]{elegantbook}

% ---------- 可复现的中文字体 ----------
\setCJKmainfont[
  BoldFont={FandolSong-Bold.otf},
  ItalicFont={FandolKai-Regular.otf}
]{FandolSong-Regular.otf}
\setCJKsansfont[
  BoldFont={FandolHei-Bold.otf}
]{FandolHei-Regular.otf}
\setCJKmonofont{FandolFang-Regular.otf}

\setCJKfamilyfont{zhsong}[
  BoldFont={FandolSong-Bold.otf}
]{FandolSong-Regular.otf}
\setCJKfamilyfont{zhhei}[
  BoldFont={FandolHei-Bold.otf}
]{FandolHei-Regular.otf}
\setCJKfamilyfont{zhkai}{FandolKai-Regular.otf}
\setCJKfamilyfont{zhfs}{FandolFang-Regular.otf}

\newcommand*{\songti}{\CJKfamily{zhsong}}
\newcommand*{\heiti}{\CJKfamily{zhhei}}
\newcommand*{\kaishu}{\CJKfamily{zhkai}}
\newcommand*{\fangsong}{\CJKfamily{zhfs}}
\renewcommand{\cbfseries}{\heiti}
\renewcommand{\citshape}{\kaishu}
\renewcommand{\cnormal}{\kaishu}
\renewcommand{\cfs}{\fangsong}

% ---------- 常用排版包 ----------
\usepackage{booktabs}
\usepackage{array}
\usepackage{mathtools}
\usepackage{microtype}
\usepackage{tasks}

% ---------- ElegantBook Blue 主题的固定四色 ----------
% 章节、页眉页脚、普通题目框：沉稳蓝
\definecolor{structurecolor}{RGB}{60,113,183}
% 定义框：绿色
\definecolor{main}{RGB}{0,166,82}
% 定理、引理、推论以及“证明”标签：橙色
\definecolor{second}{RGB}{255,134,24}
% 命题框、知识补充框：亮蓝色
\definecolor{third}{RGB}{0,174,247}

% ---------- 页面节奏 ----------
\setlength{\parskip}{0.25em}
\setlength{\emergencystretch}{2em}
\raggedbottom
\allowdisplaybreaks[2]



\tcbset{
  common/.style={
    fontupper=\citshape,
    lower separated=false,
    coltitle=white,
    colback=gray!5,
    boxrule=0.5pt,
    fonttitle=\bfseries,
    enhanced,
    breakable,
    top=8pt,
    before skip=8pt,
    attach boxed title to top left={
      yshift=-0.11in,
      xshift=0.15in
    },
    boxed title style={
      boxrule=0pt,
      colframe=white,
      arc=0pt,
      outer arc=0pt
    },
    separator sign={.}
  },
  defstyle/.style={
    common,
    colframe=main,
    colback=main!5,
    colbacktitle=main,
    overlay unbroken and last={
      \node[anchor=south east,outer sep=0pt]
        at (\linewidth-width,0)
        {\textcolor{main}{$\clubsuit$}};
    }
  },
  thmstyle/.style={
    common,
    colframe=second,
    colback=second!5,
    colbacktitle=second,
    overlay unbroken and last={
      \node[anchor=south east,outer sep=0pt]
        at (\linewidth-width,0)
        {\textcolor{second}{$\heartsuit$}};
    }
  },
  prostyle/.style={
    common,
    colframe=third,
    colback=third!5,
    colbacktitle=third,
    overlay unbroken and last={
      \node[anchor=south east,outer sep=0pt]
        at (\linewidth-width,0)
        {\textcolor{third}{$\spadesuit$}};
    }
  }
}

% 证明位于定理框外：橙色粗体“证明”标签，正文为近黑色仿宋。
% ElegantBook v4.7 默认不在结尾添加 QED 方块。
\renewenvironment{proof}{
  \par\noindent
  \textbf{\color{second}\proofname\;}
  \color{black!90}\cfs
}{
  \par
}

\newtcolorbox{checkedproblem}[1]{
  enhanced,
  colback=white,
  colframe=structurecolor!82,
  colbacktitle=structurecolor,
  coltitle=white,
  fonttitle=\bfseries,
  title={#1},
  boxrule=0.7pt,
  arc=1.5mm,
  left=7pt,
  right=7pt,
  top=8pt,
  bottom=7pt,
  before skip=10pt,
  after skip=10pt,
  attach boxed title to top left={
    xshift=5mm,
    yshift=-2.5mm
  },
  boxed title style={
    arc=1mm,
    boxrule=0pt
  }
}

\newtcolorbox{supplementbox}[1][知识补充]{
  enhanced,
  breakable,
  colback=third!4,
  colframe=third!75,
  title={#1},
  colbacktitle=third!85,
  coltitle=white,
  fonttitle=\bfseries,
  boxrule=0.6pt,
  arc=1.2mm,
  left=7pt,
  right=7pt,
  top=7pt,
  bottom=7pt
}

\newtcolorbox{analysisbox}[1][解析]{
  enhanced,
  breakable,
  colback=structurecolor!3,
  colframe=structurecolor!58,
  title={#1},
  colbacktitle=structurecolor!88,
  coltitle=white,
  fonttitle=\bfseries,
  boxrule=0.6pt,
  arc=1.2mm,
  left=7pt,
  right=7pt,
  top=7pt,
  bottom=7pt,
  before skip=6pt,
  after skip=12pt
}

\newlength{\problemSlotHeight}

% 可选参数用于在该页顶部生成一个 section 标题。
\newenvironment{twoproblempage}[1][0.46\textheight]{%
  \clearpage
  \begingroup
  \setlength{\parskip}{0pt}%
  \setlength{\problemSlotHeight}{#1}%
}{%
  \par
  \endgroup
  \clearpage
}

\newenvironment{halfpageproblem}[1]{%
  \noindent
  \begin{minipage}[t][\problemSlotHeight][t]{\textwidth}
  \begin{checkedproblem}{#1}
}{%
  \end{checkedproblem}
  \vfill
  \end{minipage}
  \par
  \vspace{0.02\textheight}
}

\newcommand{\answerspace}[1][6cm]{\par\vspace*{#1}}
